% !PRINTED 644
\weberpage{644}

Wie in § 175 beweist man die Sätze:

7. \gesperrt{Summe, Differenz, Produkt von ganzen Funktionalen in $\overline{\Omega}$ sind wieder ganze Funktionale.}

8. \gesperrt{Ist $\overline{\omega}$ ein ganzes Funktional in $\overline{\Omega}$ nach der Definition 6. und zugleich rational, so ist es ein ganzes Funktional in $\overline{Z}$ (nach der Definition 2.).}

9. \gesperrt{Ist $\overline{\omega}$ ein ganzes Funktional in $\overline{\Omega}$, so ist $N(\overline{\omega})$ ein ganzes Funktional in $\overline{Z}$, und die Absolute von $N(\overline{\omega})$ heißt die \glqq absolute Norm von $\overline{\omega}$\grqq{}.}

Die absolute Norm von $\overline{\omega}$ ist also eine Funktion in $Z$.

Bezeichnen wir die absolute Norm mit $N_a(\overline{\omega})$, so gilt der Satz:
\weq{2}{N_a(\overline{\omega}\,\overline{\omega}') = N_a(\overline{\omega})\,N_a(\overline{\omega}').}

\begin{center}
§ 180. \gesperrt{Teilbarkeit von Funktionalen. Einheiten.}
\end{center}

Es sind nun die Sätze von Bd. II, § 155 mit ganz geringen Modifikationen, auf die im folgenden aufmerksam gemacht ist, zu wiederholen. Zur Vereinfachung des Ausdruckes sollen hier mit den kleinen griechischen Buchstaben ganze Funktionale in $\overline{\Omega}$, mit den kleinen lateinischen Buchstaben ganze Funktionale in $\overline{Z}$ bezeichnet sein. Dann haben wir:

1. \gesperrt{Definition: Wenn $\beta$ von Null verschieden ist, so heißt $\alpha$ durch $\beta$ teilbar, wenn $\alpha/\beta = \gamma$ ein ganzes Funktional in $\overline{\Omega}$ ist.}

2. \gesperrt{Sind $\xi, \eta, \ldots$ beliebige ganze Funktionale in $\overline{\Omega}$ und $\alpha, \beta, \ldots$ durch $\delta$ teilbar, so ist auch}
\[
\xi\alpha + \eta\beta + \cdots
\]
\gesperrt{durch $\delta$ teilbar.}

3. \gesperrt{Definition: Ein ganzes Funktional $\varepsilon$, dessen Reziprokes $1/\varepsilon$ ganz ist, d.\,h. ein Teiler der Zahl 1 heißt eine Einheit in $\overline{\Omega}$. Eine Einheit ist Teiler eines jeden ganzen Funktionals. Produkt und Quotient zweier Einheiten sind wieder Einheiten.}

4. \gesperrt{Zwei ganze Funktionale $\alpha$, $\beta$, die gegenseitig durcheinander teilbar sind, heißen assoziiert. Ihr Quotient $\alpha/\beta = \varepsilon$ ist eine Einheit. Zwei Funktionale $\alpha$ und $\alpha\varepsilon$ sind assoziiert.}

5. \gesperrt{Ist $\alpha$ teilbar durch $\beta$, so ist jedes mit $\alpha$ assoziierte Funktional teilbar durch jedes mit $\beta$ assoziierte Funktional.}
